% BEGIN PROFESSOR FEEDBACK PF26: Discussion - Temporal window
\subsection{Choice of the temporal window}

RQ2 asks how window size affects fingerprint usefulness. Over the tested range,
longer windows generally increase both detection AP and the number of addresses
with a nonzero fingerprint. This answers the question without requiring a unique
optimum. The
40-role HuggingFace analysis selects 30 days for both cohorts, while two weeks
remains within one standard error of the maximum. Both 32-role controlled
cohorts select two weeks, and all three prevalence cohorts select 30 days. The
Real-CATS binary
comparison also selects 30 days, but a window of only three days is within one
standard error of its maximum. These results point towards a broader range of
longer windows rather than one clearly best setting.

Motif coverage helps explain this pattern. In the 40-role HuggingFace analysis,
coverage increases from 12.2\% to 77.7\% in the controlled cohort and from
11.1\% to 75.9\% in the prevalence cohort between five minutes and 30 days.
A short window leaves most
addresses with an all-zero fingerprint because their relevant transfers do not
occur close enough together. Increasing $\delta$ allows more three-edge
sequences to be counted and gives the model information for more addresses. The
parallel increase in coverage and AP suggests that part of the benefit of a
longer window comes from reduced sparsity. The experiment does not show how much
of the improvement is caused by coverage and how much comes from a more suitable
behavioural time scale.

This distinction is important for interpreting RQ2. A selected window defines
the maximum time span of one three-edge motif. It does not measure the duration
of a fraud campaign. A 30-day window can contain a rapid transaction sequence in
the same way as a five-minute window, but it also admits sequences whose edges
are much farther apart. The selected value therefore reflects a compromise
between observing enough motifs and keeping the events temporally close. Since
30 days is the upper boundary of the tested grid, the results also cannot show
whether performance would continue to increase at still larger windows. The
supported conclusion is limited to the tested range.
% END PROFESSOR FEEDBACK PF26

